% Transcription — fonds Alexandre Grothendieck, shelfmark 29, batch 4 (pages 61-80).
% Established from the facsimile archives/batches/29/batch-04.pdf (a mirror of
% https://grothendieck.umontpellier.fr/29.pdf), per the skill
% transcribe-grothendieck.
%
% All twenty pages carry mathematics or mathematical correspondence; none is
% skipped.
%
% The batch falls into four runs.
%
% Pages 61-64 finish the run "Th. de finitude relatif pour pi_1" begun at page
% 46 of batch 3. Page 60 was Grothendieck's own page 1; pages 61, 62 and 63
% carry his 2, 3 and 4 in the top-left corner. They are the hardest pages of
% the batch: 61 is a palimpsest of pencil under ink, reworked and struck
% throughout, and much of it is left as \ill{}. What survives is the Bertini
% reduction to a curve — X' the normalisation of X_red, U = X - D, and the
% passage from the triple (k, X, U) to (K, C, V) — and the two corollaries of
% page 62-63: pi_1^t of a product of pairs is the product of the pi_1^t, and
% pi_1'(X x Y) = pi_1'(X) x pi_1'(Y). Page 64 is a scratch sheet of isolated
% fragments and diagrams on the same question, over heavy bleed-through from
% its verso.
%
% Pages 65-69 are the run Grothendieck headed "pi_1 birationnel" in the margin
% of page 65 and paginated himself A to E: the theorem that pi_1(X) is
% topologically of finite type for X whose irreducible components are the
% complement of a codimension-2 closed set in a proper scheme, Zariski's
% Bertini lemma and its corollary on the surjectivity of pi_1(X_u) -> pi_1(X),
% the reduction to a generic curve of genus g, and then the passage to the
% limit over the proper normal models of a function field, defining what he
% calls the birational fundamental group pi_1(K/k, Omega) and stating its
% finite generation.
%
% Pages 70-76 are the parallel run "Pic birationnel", paginated F to L: the
% theorem that Pic^tau_Y -> Pic^tau_X is an isomorphism for a dominant
% birational morphism of integral normal proper schemes, its proof by way of
% the open U where f^{-1} is defined, the corollary on H^1(-, O), the
% characteristic-p argument through Frobenius and its kernel, the
% characteristic-0 argument by spreading out over a scheme of finite type over
% Z, and finally the uniform bound dim Pic_{X/k} <= n(K/k) over all proper
% normal models, with n(K/k) the genus of a generic curve section.
%
% Page 77 is the sheet Grothendieck headed "Revêtements formels étales et th.
% de Mumford", above a struck box; pages 78-80 are what it covers — Murre's
% letter of 30 April 1969 sending the preliminary version of those notes, and
% his two typed pages of comments on its chapters I to X. Both are in English
% and stay as written, as in batches 1 and 2.
%
% The dating is the inventory's own, for the whole shelfmark.
%
% Pass: Opus 5 (claude-opus-5), 2026-08-22 — first pass, unchecked against the pages by a human.
\documentclass[11pt,a4paper]{article}
% Le préambule commun à toutes les transcriptions du fonds Grothendieck.
%
% Il tient en une idée : **l'incertitude se compose, elle ne se résout pas.**
% Une transcription qui lisse un mot douteux a détruit la seule chose pour
% laquelle on la faisait — un lecteur doit pouvoir savoir, à chaque endroit,
% ce qui était sur le feuillet et ce qui a été ajouté. Les six macros qui
% suivent sont donc l'appareil critique, et non de la décoration.
%
% Les mêmes commandes sont reconnues par `scripts/render.mjs`, qui produit la
% vue de lecture du site. Une macro ajoutée ici doit l'être aussi là-bas,
% faute de quoi le rendu échouera bruyamment — ce qui est voulu : un rendu
% silencieusement faux est pire qu'un rendu absent.

\ProvidesPackage{grothendieck}[2026/08/08 Transcriptions du fonds Grothendieck]

\usepackage{amsmath,amssymb,amsthm}
\usepackage{mathrsfs}
\usepackage{stmaryrd}
% Les diagrammes commutatifs sont partout dans ce fonds : c'est la langue
% même de la géométrie algébrique de Grothendieck, et une transcription qui
% les rendrait en prose ne transcrirait rien.
\usepackage{tikz-cd}
% Français par défaut — c'est la langue des feuillets — mais l'anglais est
% chargé aussi : la traduction et le résumé sont des documents anglais, et la
% typographie française (espaces avant les deux-points) y serait une faute.
% Ces documents basculent par \selectlanguage{english} en tête de corps.
\usepackage[english,french]{babel}
\usepackage{fontspec}
\usepackage{geometry}
\geometry{a4paper,margin=25mm,marginparwidth=28mm}
\usepackage{marginnote}
\usepackage{xcolor}
% ulem, not soul. `soul`'s \st and \ul are built on a character-by-character
% scan that XeLaTeX's font machinery breaks: the document fails with an
% undefined control sequence rather than degrading. ulem does the same two jobs
% and is Unicode-engine safe. `normalem` stops it hijacking \emph, which the
% transcriptions use for Grothendieck's own emphasis.
\usepackage[normalem]{ulem}
\usepackage{hyperref}
\hypersetup{colorlinks=true,linkcolor=black,urlcolor=black,pdfborder={0 0 0}}

% --- Métadonnées du lot -----------------------------------------------------
% Reprises telles quelles par le script de rendu, qui les affiche en tête de
% la vue de lecture. Elles fixent ce que le fichier prétend couvrir : sans
% elles, une transcription flotte sans point d'attache dans un fonds de seize
% mille feuillets.

\newcommand{\folder}[1]{\gdef\grothFolder{#1}}
\newcommand{\batch}[1]{\gdef\grothBatch{#1}}
\newcommand{\leaves}[2]{\gdef\grothFirst{#1}\gdef\grothLast{#2}}
\newcommand{\foldertitle}[1]{\gdef\grothFoldertitle{#1}}
\gdef\grothFolder{?}\gdef\grothBatch{?}\gdef\grothFirst{?}\gdef\grothLast{?}\gdef\grothFoldertitle{}

% --- Le résumé --------------------------------------------------------------
%
% Il ouvre la lecture modernisée, sous le titre « Résumé », et il est composé
% en petit. Ce n'est pas un ornement : c'est le seul passage du document qui
% ne soit pas des mathématiques, et un lecteur doit pouvoir le reconnaître
% avant de l'avoir lu — puis le sauter s'il connaît déjà le sujet, ou s'y
% arrêter s'il ne le connaît pas. Le filet qui le suit marque la reprise du
% corps.

\newenvironment{resume}
  {\small\setlength{\parskip}{0.45em}}
  {\par\medskip\hrule\medskip}

% --- Mots-clés ---------------------------------------------------------------
%
% En anglais, en clôture du résumé de la lecture modernisée : le vocabulaire
% moderne sous lequel chercher ce que les feuillets construisent. Le manifeste
% du site (`npm run manifest`) les extrait pour étiqueter la cote entière —
% cette ligne est la source unique des tags, il n'y a pas de fichier à côté.
\newcommand{\keywords}[1]{\par\smallskip\noindent
  {\footnotesize\textsc{Keywords} --- \textit{#1}}\par}

% --- L'appareil critique ----------------------------------------------------

% Le numéro de feuillet, en marge. C'est l'ancre qui permet de régler un
% désaccord sur une formule en regardant *un* feuillet plutôt qu'une chemise
% de six cents. Le site s'en sert aussi pour tourner les pages du fac-similé
% quand on fait défiler la transcription.
\newcommand{\leaf}[1]{%
  \marginnote{\footnotesize\color{gray!70}\textsf{#1}}%
  \label{leaf:#1}\ignorespaces}

% Illisible : on ne devine pas. Un mot inventé qui se lit comme les autres est
% le pire résultat possible.
\newcommand{\ill}{\textcolor{red!60!black}{[\,\ldots\,]}}

% Lecture douteuse : le mot est donné, et signalé comme incertain.
\newcommand{\uncertain}[1]{\uline{#1}}

% Ajout de l'éditeur — une lettre manquante, un mot sous-entendu.
\newcommand{\add}[1]{\textcolor{blue!50!black}{[#1]}}

% Biffé par Grothendieck lui-même. Ce qu'il a rayé fait partie du document :
% une rature dit souvent où la pensée a changé de direction.
\newcommand{\struck}[1]{\sout{#1}}

% Note du transcripteur, sur l'état matériel du feuillet.
\newcommand{\note}[1]{\par\smallskip{\small\color{gray!60!black}\textit{[#1]}}\par\smallskip}

% Note marginale de Grothendieck, distincte du corps du texte.
\newcommand{\marginal}[1]{\par\smallskip\noindent\hspace{1em}%
  {\small\color{gray!60!black}#1}\par\smallskip}

% --- Titre ------------------------------------------------------------------

\AtBeginDocument{%
  \begin{center}
    {\large\bfseries Fonds Alexandre Grothendieck --- cote n\textsuperscript{o}~\grothFolder}\\[2pt]
    {\itshape\grothFoldertitle}\\[4pt]
    {\small feuillets \grothFirst--\grothLast\ (lot \grothBatch)}\\[2pt]
    {\footnotesize\color{gray!60!black}Transcription établie d'après le fac-similé numérisé
     par l'Université de Montpellier.\\
     Les lectures incertaines sont soulignées, les illisibles notées [\,\ldots\,],
     les ajouts entre crochets.}
  \end{center}
  \vspace{1em}}

\endinput


\folder{29}
\batch{4}
\pages{61}{80}
\dating{1959-1969}
\watermark{Édition de démonstration}
\foldertitle{Groupe fondamental [Autour de SGA 4 et SGA 7] : lettres (1959, 1967, 1969), tapuscrits et copies de tapuscrit annotés (s.d.), notes manuscrites (s.d.)}

\begin{document}

\section*{« Th.\ de finitude relatif pour $\pi_1$ », fin (pages 61 à 64)}

\page{61}
\note{Grothendieck numérote cette page 2. La page est un palimpseste : un
premier jet au crayon, repris à l'encre, puis biffé par endroits. La plus
grande partie n'est pas lisible à la résolution du fac-similé et reste
\ill{}.}

\emph{Théorème.} \struck{\ill{}} $X$, \struck{préschéma} \ill{} alg.\ sur
\ill{} clos, \ill{}

\ill{} $X'$ le normalisé de $X_{\mathrm{red}}$, $D_i$ les composantes
\add{irréductibles} de codim.\ 1 de $p^{-1}(D)$ \ill{}

\begin{tikzcd}
X' \arrow[d, "p"] & \supset & D'_i \arrow[d] \\
X & \supset & D
\end{tikzcd}

$$U = X - D$$

\ill{} $R^0 p'_{*} \ill{}$ \ill{} si $Y$ est régulier, il suffit \ill{}

\ill{}

Puis, en bas de page, la réduction à une courbe :

\begin{tikzcd}
C \arrow[r] \arrow[d, hook'] & X \arrow[d, hook'] \\
V \arrow[r] & U
\end{tikzcd}

\begin{tikzcd}
C \times Y = C_Y \arrow[r] \arrow[d, hook'] & R \arrow[d, hook'] \\
V \times Y = V_Y \arrow[r] & P
\end{tikzcd}

\ill{} $\pi_1(C) \to \pi_1(X)$ \struck{\ill{}} \ill{} les \ill{} images des
$D_i$ dans $C$ sont des diviseurs \ill{} \uncertain{singuliers}, \ill{} \ill{}
où \ill{} l'induction \ill{} \ill{} sur $(k, X, U)$ \ill{} par $(K, C, V)$,
\ill{} d'une \emph{courbe}. Mais alors, \ill{} \ill{} les conditions, \ill{}

\page{62}
\note{Grothendieck numérote cette page 3}

\struck{d'application \ill{} la théorie de spécialisation \ill{} des \ill{}
\ill{}}

Lorsque $X$ n'est pas normal, on utilise le \uncertain{dévissage} \ill{}
À détailler.

\emph{Corollaire.} \struck{\ill{}} Soit $U$ de type fini sur $k$, $Y$
\struck{\ill{}} \add{de type fini sur $k$}, $R$ un revêtement étale
\add{soit} de $U \times_k Y$, \ill{} $Y$ \ill{} Supposons que pour \ill{}
$y$, $R_y$ est un revêtement \ill{} trivial. Alors $R$ \ill{}, \ill{}
au-dessus d'un voisinage de \struck{\ill{}} $y$, \ill{} \ill{} d'un \ill{}
\ill{} d'un tel voisinage.

\emph{Cor.} Soient $X'$, $X''$ propres sur $k$, $Z'$ ($Z''$) des parties
fermées, \ill{} $X = X' \times X''$, $Z = Z' \times X'' \cup X' \times Z''$,
\ill{} considérant les groupes $\pi_1^t(X'; Z'; x')$, $\pi_1^t(X''; Z''; x'')$
et $\pi_1^t(X; Z, x)$, \struck{\ill{}} notés $\pi$, $\pi'$, $\pi''$. Alors

\[ \struck{\pi_1^t(X)} \quad \pi \longrightarrow \pi' \times \pi'' \]

est un \emph{isomorphisme}.

\emph{Cor.} Soient $X$, $Y$ de type fini sur $k$, \ill{}
$\pi_1'(X)$ et $\pi_1'(Y)$, \add{$\pi_1'(X \times Y)$} [mais elles \ill{}
\ill{} la car.]. Alors

\[ \pi_1'(X \times Y) \;\xrightarrow{\ \sim\ }\; \pi_1'(X) \times \pi_1'(Y). \]

\page{63}
\note{Grothendieck numérote cette page 4}

\begin{tikzcd}
X' \arrow[d] \\
X
\end{tikzcd}

$X'$ un normalisé ; \quad $Y \subset X$ --- \uncertain{complète},
\uncertain{irréductible} ; \quad $U = X - Y$.

$Y_1 \ldots Y_r$ les composantes de $Y$ de codim.\ 1 dans $X$ ;
$Y'_{ij}$ les images inverses des $Y_i$ dans $X'$.

\struck{$\pi_1(U_i)$} \quad \struck{$\pi_1(U)$ le quotient de $\pi_1(U)$ par
le noyau de $\pi_1(U) \to \pi_1(X)$} \quad \struck{est \ill{} engendré par}

Considérons les rev.\ étales de $X - Y$ qui \struck{sont classifiés}
deviennent étales sur les \uncertain{données} des $\mathcal{O}_{Y'_{ij}}$ :
ils sont classifiés par un groupe de type galoisien $\pi$.

\emph{Th.} Le groupe $\pi$ est top.\ de génération finie.

\emph{Cor.\ 1.} Soit $X$ \ill{} complète, $Y$ de \ill{} $\geq 2$,
\add{$U = X - Y$}. Alors $\pi_1(U)/\ill{}$ est top.\ de génération finie.

\emph{Cor.\ 2.} Soit $X$ \struck{de type} alg. Alors $\pi_1'(X)$ est top.\ de
type fini.

\page{64}
\note{Feuille de brouillon : des fragments isolés, sans phrase liante, sur la
même question. Ils sont transcrits dans l'ordre où ils viennent sur la page.}

\[ \pi_1(X \times Y) \;\xrightarrow{\ \mathrm{surj}\ }\;
   \pi_1(X) \times \pi_1(Y) \]

$X \supset Z$ \qquad $\pi_1(Z) \to \pi_1(X) \to \pi_1(Y)$

\begin{tikzcd}
X \arrow[d] \\
Y
\end{tikzcd}
\qquad
\begin{tikzcd}
X'' \arrow[d, Rightarrow] \\
X' \arrow[d] \\
X \arrow[d] \\
{}
\end{tikzcd}

\begin{tikzcd}
Y' \arrow[r] \arrow[d] & X' \arrow[d] \\
Y \arrow[r, hook] \arrow[d] & X \arrow[d] \\
{} & \mathbf{P}^r
\end{tikzcd}
\qquad
\begin{tikzcd}
X' \arrow[r, hook] \arrow[d] & \bar X' \arrow[d] \\
X \arrow[r, hook] \arrow[d] & \bar X \arrow[dl] \\
\mathbf{P}^r
\end{tikzcd}
\qquad
\begin{tikzcd}
Y \arrow[r] \arrow[d] & X \arrow[d] \\
H \arrow[r] & \mathbf{P}^r
\end{tikzcd}

\marginal{surj ??}

\[ \pi_1(Y) \longrightarrow \pi_1(X) \qquad (X \text{ normal}) \]

$X \subset \hat X$ \ill{} ; \quad $X_0 \subset \hat X_0$ \uncertain{plus
régulier}

\[ \pi_1(X) \qquad\qquad
   \struck{\pi_1(Y) \to \pi_1(X) \leftarrow \pi_1(X - Y)} \]

\section*{« $\pi_1$ birationnel » (pages 65 à 69)}

\page{65}
\note{Grothendieck pagine lui-même cette suite A à L ; la page 65 porte A.}

\marginal{$\pi_1$ birationnel}

\emph{Th.} Soient $k$ un corps alg.\ clos, $X$ schéma \add{connexe}
\struck{propre} sur $k$, \struck{$Y$ une partie fermée} tel que les
composantes irréductibles $X_i$ soi\add{en}t isomorphes (avec la structure
réduite \uncertain{induite}) au \struck{la complémentaire} complémentaire
d'une partie fermée de codim.\ $\geq 2$ dans un schéma propre sur $k$.

\marginal{Question. $Y \subset X$ fermé, de codim.\ $\geq 2$ : est-ce que
$\pi_1(X-Y) \to \pi_1(X)$ \uncertain{a un noyau} \ill{} de type fini ??}

Alors $\pi_1(X)$ est top.\ de type fini.

\emph{Démonstration.} Par SGA IX 6.2, on est ramené au cas où $X$ est
\struck{irréductible} \add{intègre}, donc de la forme $X = Z - Y$, où $Z$ est
propre et $Y$ de codim.\ $\geq 2$ dans $Z$. \ill{} \ill{}

\[ X = \bar X - \bar X \cap Y, \]

et $\bar X$ \uncertain{étant} une composante irréductible de $Z$,
$\bar X \cap Y$ qui est de codim.\ $\geq 2$. \ill{} On peut supposer $Z$
\struck{réduit} intègre, et on est ramené à prouver le

\emph{Corollaire.} $Z$ intègre et propre sur $k$, $Y$ de codim.\ $\geq 2$
dans $Z$, $X = Z - Y$ ; alors $\pi_1(X)$ est top.\ de type fini.

Soit $Z' \to Z$ un morphisme propre surjectif, avec $Z'$ intègre ; soit $X'$
l'image inverse de $X$ dans $Z'$, soit $X' = Z' - Y'$, et $Y'$ de codim.\
$\geq 2$ \add{dans $Z'$}. \ill{} il suffit de prouver que $\pi_1(X')$ est
top.\ de type fini. \uncertain{Par} le lemme de Chow,

\page{66}
on peut donc supposer que $Z$ est projectif, \struck{si dim} et \ill{} par
récurrence. \add{$\dim Z \leq 0$, trivial.} Si $\dim Z = 1$, on aura
$Y = \emptyset$, i.e.\ $X = Z$, et le th.\ est connu. Si $\dim Z \geq 2$ on
utilise ceci :

\emph{Lemme de Bertini} (Zariski). Soit $k$ un corps, $X \xrightarrow{f}
\mathbf{P}^r_k$ un morphisme, avec $X$ irréductible, $\dim f(X) \geq n$ ;
soit $1 \leq s \leq n-1$, et soit $u$ le point générique de
$\mathbf{P}^r_k \times \cdots \times \mathbf{P}^r_k$ ($s$ facteurs), et $K$
une clôture algébrique de $k(u)$, enfin $L_u$ la variété linéaire dans
$\mathbf{P}^r_K$ définie par $u$, et

\[ X_u = X \times_{\mathbf{P}^r_k} L_u
       = X_K \times_{\mathbf{P}^r_K} L_u. \]

\marginal{peut-on remplacer « irréductible » par « connexe » ? --- non}

Alors $X_u$ est irréductible.

\emph{Corollaire.} \struck{Il est} Si $X$ est normal, alors
$\pi_1(X_u) \to \pi_1(X)$ est surjectif.

\begin{tikzcd}[column sep=small, row sep=small, nodes={font=\scriptsize}]
X'_u \arrow[r] \arrow[d] & X' \arrow[d] \\
X_u \arrow[r] \arrow[d] & X_K \arrow[r] \arrow[d] & X \arrow[d] \\
L_u \arrow[r] & \mathbf{P}^r_K \arrow[r] & \mathbf{P}^r_k
\end{tikzcd}

En effet, si $X' \to X$ est un revêtement \struck{connexe} étale connexe de
$X$, \ill{} on montre \ill{} irréductible, donc
$X'_u = X' \times_X X_u$ est irréductible, donc connexe, O.K.

\marginal{\ill{}}

\page{67}
Reprenons donc $Z$ normal de dim.\ $n \geq 2$, et $X = Z - Y$. Plongeons $Z$
dans $\mathbf{P}^r$, et le morphisme induit $f : X \to \mathbf{P}^r$. On
prend $s = n-1$, de sorte que
$X_u = X_K \cap L_u = Z_K \cap L_u$ \struck{est une courbe} $\neq \emptyset$,
car $Y_K \cap L_u = \emptyset$ ; donc $X_u = Z_K \cap L_u$ est une courbe
\struck{\ill{}} projective, d'ailleurs normale. Si $g$ est le genre, il
résulte alors du corollaire au lemme de Bertini que $\pi_1(X)$ admet $2g$
générateurs topologiques, O.K.

\emph{Conséquence.} Soit $K$ une extension de type fini d'un corps $k$. Je
\ill{} existe toujours un modèle projectif, \struck{celui-ci} \struck{normal}
\ill{} propre, et \ill{} \struck{\ill{}}

\struck{\ill{} $X'$ \ill{} les modèles \ill{} particuliers \ill{} régulier
$U$ \ill{} « presque propre », i.e.\ complémentaire d'un \ill{} fermé de
codim.\ $\geq 2$ dans \ill{} \ill{} une partie fermée de codim.\ $\geq 2$
dans \ill{} que \ill{} modèle propre $X$.}

\begin{tikzcd}
& X'' \arrow[dl] \arrow[dr] & \\
X & & X'
\end{tikzcd}

\struck{D'ailleurs, tous les modèles $U$ comme dessus, ordonnés par \ill{}}

\page{68}
D'ailleurs les modèles propres et normaux, pour la relation de domination,
forment un ordonné filtrant croissant. \struck{\ill{}} Soit $X_0$ un tel
modèle \ill{} \add{propre normal} \ill{} régulier, et sur $X_0$ \ill{}
\struck{\ill{}}

\begin{tikzcd}
X_1 \arrow[d, "f"] & X_1 \arrow[l] \\
X_0 & \supset U_1 \supset U_0
\end{tikzcd}

$U_1 \to X_1$ induit \struck{\ill{}} \ill{} et à la limite, c'est \ill{} de
codim.\ $\geq 2$, \struck{\ill{} complémentaire dans $X_0$ d'un \ill{} fermé
de codim.\ $\geq 2$. \ill{}}

\begin{tikzcd}
\pi_1(X_1) \arrow[d, "f_{*}"'] & \pi_1(U_1) \arrow[l] & \\
\pi_1(X_0) & \pi_1(U_0) \arrow[l] \arrow[r, "\sim"] \arrow[ul, "g_{*}"] & \pi_1(U_1)
\end{tikzcd}

\note{sous l'isomorphisme du bas il écrit « th.\ \add{de} pureté »}

L'injection \ill{} induite $U \to X$ étant \ill{} le complémentaire d'une
partie fermée de codim.\ $\geq 1$, on voit que (\ill{} de pureté) on a un
homo.\ \emph{surjectif} (si $X$ normal)

\[ \pi_1(U) \longrightarrow \pi_1(X) \]

d'où : à la limite un homo.\ surjectif

\[ \pi_1(U, \Omega) \longrightarrow
   \varprojlim_{X \text{ modèle propre}} \pi_1(X, \Omega) \]

où $\Omega$ est une extension alg.\ close de $K$.

\page{69}
Le dernier est ce que j'ai appelé le \emph{groupe fondamental birationnel} de
$K/k$, \ill{} noté $\pi_1(K/k, \Omega)$. \ill{} tout modèle \ill{}
$\pi_1(X, \Omega)$ \ill{} quotient de $\pi_1(U)$. \ill{} Il \ill{} être
\ill{} complémentaire d'une partie fermée de codim.\ $\geq 2$ dans \ill{}

\marginal{\ill{}}

\emph{Th.} Si $K/k$ est une extension de type fini, avec $k$ alg.\ clos,
alors \struck{\ill{}} $\pi_1(K/k, \Omega)$ est top.\ de type fini.

\section*{« Pic birationnel » (pages 70 à 76)}

\page{70}
\marginal{Pic birationnel}

\emph{Théorème.} Soient $f : X \to Y$ un morphisme dominant
\add{birationnel} de préschémas intègres normaux, \struck{et} propres sur
$k$, avec $Y$ \struck{\ill{}} \add{loc.\ régulier}. Alors le morphisme

\[ \underline{\mathrm{Pic}}^{\tau}_Y \longrightarrow
   \underline{\mathrm{Pic}}^{\tau}_X \]

est un \emph{isomorphisme}.

\marginal{$X$ normal ? « sur $k$ » ? --- \ill{} $Y$ régulier ?? \ill{} sur
$k$ ? --- N.B.\ je \ill{}}

C'est un \emph{monomorphisme} grâce au fait que
$f_{*}(\mathcal{O}_X) = \mathcal{O}_Y$ (universellement sur $k$ !), donc une
immersion fermée puisque un homo.\ de groupes \ill{} \ill{} propre. Montrons
que c'est \emph{surjectif}. Soit $U$ l'ouvert de $Y$ où $f^{-1}$ est défini ;
comme $Y$ est factoriel, \ill{}

\[ \mathrm{Pic}(X) \xrightarrow{\ \sim\ } \mathrm{Pic}(U), \]

d'où, utilisant

\begin{tikzcd}
& X \arrow[d, "f"] \\
U \arrow[ur, "g"] \arrow[r, "i"'] & Y
\end{tikzcd}

le diagramme commutatif

\begin{tikzcd}
& \mathrm{Pic}(X) \arrow[dl, "g^{*}"] \\
\mathrm{Pic}(U) & \mathrm{Pic}(Y) \arrow[u, "f^{*}"'] \arrow[l, "i^{*}"]
\end{tikzcd}

qui montre que $\mathrm{Pic}(Y) \to \mathrm{Pic}(X)$ est un homo.\ \ill{} un
\emph{facteur direct}, l'image étant \ill{} $\xi$ tels que
\struck{$f^{*}(\xi) = 0$} $\xi = f^{*} i^{*-1} g^{*}(\xi)$. \ill{} $\xi$
\ill{} $\tau$-équivalent à 0, donc $i^{*-1} g^{*}(\xi)$ aussi,
\struck{\ill{}} \ill{} $= \ill{}\ f^{*} i^{*-1} g^{*}(\xi)$. D'autre part, on
voit \ill{}

\marginal{Revoir tout ça avec les cycles !}

\page{71}
diviseur $\sum n_i X_i$, où les $X_i$ sont les « diviseurs
\uncertain{singuliers} » pour $f$. Il reste donc à voir que si une telle
combinaison linéaire est \struck{algébriquement} $\tau$-équivalente à 0, elle
est nulle. \struck{Ce} (C'est là un lemme bien connu\ldots)

\marginal{vérifier ? \ill{}}

Ainsi, $f^{*} : \underline{\mathrm{Pic}}^{\tau}_Y \to
\underline{\mathrm{Pic}}^{\tau}_X$ est une \emph{immersion surjective}. Il
reste à prouver qu'elle induit un isom.\ sur les \ill{} tangents à
l'origine, donc à vérifier le

\marginal{on peut ramener au projectif --- \textbf{insuffisant}}

\emph{Corollaire.} Sous les conditions du th., l'homomorphisme

\[ f^{*} : H^1(Y, \mathcal{O}_Y) \longrightarrow H^1(X, \mathcal{O}_X) \]

est un isomorphisme.

On sait déjà qu'il est injectif, la question est la surjectivité.
\struck{Il suffit} En caractéristique $p > 0$, il suffit de prouver que sont
bijectifs

\[ H^1(Y, \mathcal{O}_Y)^F \longrightarrow H^1(X, \mathcal{O}_X)^F \]
\[ {}_F H^1(Y, \mathcal{O}_Y) \longrightarrow {}_F H^1(X, \mathcal{O}_X) \]

où $F$ est le Frobenius. Le premier homo.\ s'écrit

\[ H^1(Y, \mathbf{Z}/p\mathbf{Z}) \longrightarrow
   H^1(X, \mathbf{Z}/p\mathbf{Z}) \]

\page{72}
c'est un isomorphisme : cela \ill{} du fait que $\pi_1(X) \to \pi_1(Y)$ est
un isomorphisme [\ill{} le \uncertain{th.\ de pureté} sur $Y$].

\marginal{$Y$ régulier}

Le deuxième homomorphisme s'écrit

\[ H^0(Y, \mathcal{O}_Y / \mathcal{O}_Y^{p^{\nu}}) \longrightarrow
   H^0(X, \mathcal{O}_X / \mathcal{O}_X^{p^{\nu}}) \]

[grâce au fait que $X$, $Y$ sont réduits], ou aussi, supposant que $k$ est
alg.\ clos donc parfait, s'interprète comme

\[ f^{*} : H^0(Y, \mathcal{O}_Y^{p^{-\nu}} / \mathcal{O}_Y) \longrightarrow
   H^0(X, \mathcal{O}_X^{p^{-\nu}} / \mathcal{O}_X) \]

\marginal{$Y$ intègre sur $k$ (?)}

Ce dernier est bijectif grâce à la normalité de $X$ et $Y$ : \ill{} \ill{}
les notations ci-dessus, on a un diagramme commutatif

\begin{tikzcd}
& T(X) \arrow[dl, "g^{*}"] \\
T(U) & T(Y) \arrow[u, "f^{*}"'] \arrow[l, "i^{*}"]
\end{tikzcd}

avec $g^{*}$ \emph{injectif} [car $\mathcal{O}_X^{p^{-\nu}} / \mathcal{O}_X$
est sans torsion \ldots], donc $f^{*}$ et \uncertain{$i^{*-1} g^{*}$} sont
\ill{} réciproques l'une de l'autre.

\marginal{$X$ normal}

\page{73}
Dans le cas où $k$ est de car.\ 0, on peut supposer que $k$ est
\struck{de type fini} \struck{sur $\mathbf{Q}$}, le corps des fonctions d'un
schéma intègre $S$ de type fini sur $\mathbf{Z}$, et que $f : X \to Y$
provient d'un morphisme de $S$-préschémas, $\bar f : \bar X \to \bar Y$,
intègres, $\bar Y$ \uncertain{lisse} sur $S$, et $\bar X$ \struck{\ill{}}
normal sur $S$, les fibres de l'un et de l'autre étant \struck{\ill{}}
connexes. Alors \struck{\ill{}}

\struck{$\forall y \in Y$, si $H^1(Y, \mathcal{O}_Y) \to
H^1(X, \mathcal{O}_X)$ \ill{} de dim.\ $r$, il existe \ill{}}

\struck{\ill{}} On peut supposer \ill{} $X$ et $Y$ alg.\ plats sur $S$
\struck{de dim.\ 1}, \struck{\ill{}} \struck{\ill{}}
$R^1 q_{*}(\mathcal{O}_{\bar Y}) \to
R^1 p_{*}(\mathcal{O}_{\bar X})$ est injectif, et \struck{que} (car il l'est
\ill{} \ill{} par \ill{} d'ailleurs) alors \ill{} \ill{} sur $S$, \ill{}
le \ill{} de $H^1(Y, \mathcal{O}_Y) \to H^1(X, \mathcal{O}_X)$ est de dim.\
$r$, \ill{} $=$ \ill{} même ; mais pour les fibres \ill{} le \emph{th.} \ill{}
les \ill{} $s \in S$, on \uncertain{vérifie} \ill{} \ill{} au \ill{} de

\page{74}
car \ill{} $> 0$. \uncertain{Cela pour} $\nu = 0$, et achève la
démonstration.

\emph{Corollaire 1.} Invariance birationnelle de
$\underline{\mathrm{Pic}}^{\tau}_{X/k}$ pour $X$ propre et simple sur $k$.

\emph{Corollaire 2.} Soit $K$ une extension \struck{régulière}
\add{séparable primaire} de type fini de $k$, admettant un modèle
\struck{régulier} propre et \emph{simple} sur $k$. Alors
$\varinjlim_X \underline{\mathrm{Pic}}^{\tau}_{X/k}$, où $X$ parcourt les
modèles propres et normaux \add{« sur $k$ »} de $K$, existe \ill{} \ill{}
[le système inductif \struck{projectif} est essentiellement constant].

\marginal{N.B.\ si $k$ \ill{}, et $K$ de dim.\ 1 sur $k$, \ill{} c'est \ill{}
\ill{} \ill{} pour idiots !}

\emph{Question.} Le cor.\ 2 reste-t-il valable en supposant seulement que $K$
admette un modèle « normal sur $k$ » ?? \struck{\ill{}} \add{le cas}
\struck{\ill{}} \uncertain{équivaudrait-il} à conjecturer une assertion,
semblables

\begin{enumerate}
\item $\dim \underline{\mathrm{Pic}}_{X/k}$ reste majoré ;
\item ordre $(\mathrm{Tors}\ \mathrm{Pic}(X/k))$ reste majoré ;
\item $\dim H^1(X, \mathcal{O}_X)$ reste majoré.
\end{enumerate}

\marginal{N.B.\ (i) résulte de (iii)}

\page{75}
\emph{Théorème.} Soit $K$ une extension de type fini du corps $k$,
\struck{ayant un \ill{} modèle propre \ill{} et universellement \ill{} sur
$k$}. \struck{Soit $(X_i)_{i \in I}$ le système \ill{}} Il existe un entier
$n(K/k)$ tel que pour tout modèle $X$ de $K$ \add{propre sur $k$, et normal
sur $k$}, \ill{} \struck{\ill{}}

\[ \dim \underline{\mathrm{Pic}}_{X/k} \leq n(K/k) \]

\emph{Démonstration.} \struck{On est ramené au cas} \add{b) il \ill{} un}
modèle $X$ de type indiqué, \ill{} traité. Sinon, $K$ est séparable sur $k$,
et

\[ \bar K = K \otimes_k \bar k \]

est une extension de type fini du corps alg.\ clos $\bar k$, et on est
ramené au cas où $k$ est alg.\ clos. Alors

\[ \dim \underline{\mathrm{Pic}}_{X/k} = \tfrac{1}{2}\,
   \mathrm{rang}_{\mathbf{Z}_{\ell}}
   \left(\text{partie } \ell\text{-primaire de } \pi_1(X/k)
   \text{ rendu abélien}\right) \]

et à ce titre, on \uncertain{se sert} du théorème de majoration uniforme du
$\pi_1(X/k)$

\page{76}
pour les modèles normaux (résultant de la partie facile de Lefschetz, en
\ill{} \ill{} les groupes fondamentaux rendus abéliens), \ill{} \ill{} fini.
De façon plus précise il suffit, prenant un \add{modèle} $X$ projectif sur
$k$ (propre et normal), de prendre dans $X$ une « courbe générique $C$ » et
pour $n(K)$ le genre de ladite courbe générique.

\begin{tikzcd}
C' \arrow[r, hook] \arrow[d] & X' \arrow[d] \\
C_{k(u)} \arrow[r, hook] & X \arrow[r, hook] & \mathbf{P}^r
\end{tikzcd}

Y a-t-il une démonstration du fait que cela donne une majoration,
n'utilisant \emph{pas} la théorie du groupe fondamental et le fait que
$\underline{\mathrm{Pic}}^{\tau}_{X/k}$ est propre sur $k$ si $X$ est
normale ?

\marginal{démonstration directe ???}

\section*{« Revêtements formels étales et th.\ de Mumford » (pages 77 à 80)}

\page{77}
\struck{Comportement d'une \ill{} kummérien\ill{} et \ill{} leur variante
($\pi_1$ discret etc.\ \ldots)}

\textbf{Revêtements formels étales et th.\ de Mumford}

\subsection*{J.-P. Murre à A. Grothendieck, Voorschoten, 30 avril 1969 (page 78)}

\page{78}
Voorschoten, 30-4-69. Palestrinalaan 11.

Dear Grothendieck,

Last fall I wrote you that I was planning to work out your notes on formal
tamely ramified coverings together with a result on ``plumbing''.

In doing so it turned out to be necessary --- at least by the method I
followed --- to write again a lot of preliminaries on tamely ramified
coverings, this time for formal schemes. I am still not ready with the above
mentioned ``project''! Since it took me already much time, I would appreciate
it very much to have --- before proceeding --- your opinion about the
\emph{preliminary version} which I have prepared so far. Of course it is not
necessary to read the notes in detail. My question is: what do you think of
this set-up and do you think it worthwhile to continue in this way? To me the
entire thing seems somewhat indigestible.

I have adjoined some comments on the different chapters. The notes consists
of two parts. Part one (Chap.\ I--VII) are the preliminaries, part two (Chap.\
VIII--X) are your notes and the plumbing. The chapter on the analogue of the
theorem of Mumford has still to be written.

I hope you don't mind too much that I bother you again with these tamely
ramified coverings!

With best wishes, Sincerely yours, J. P. Murre.

\subsection*{« Comments », deux pages dactylographiées jointes à la lettre (pages 79 et 80)}

\page{79}
Comments.

\emph{Chap.\ I.} Étale morphisms \add{of formal schemes} are defined as
inductive systems of étale morphisms of usual schemes (1.1.2 and EGA I
10.12.3). Crucial lemmas: 1.1.6, 1.1.13 and 1.4.9.

As to 1.1.13: étale nbhs for formal schemes are locally completions of étale
nbhs of usual schemes. This is used in 7.1.5, itself used in the reduction
step 9.3.3.

\emph{Chap.\ II.} There are Kummer coverings (II-1), generalized Kummer
coverings (II-2) and coverings of Kummer type (II-6)!

No attempt has been made towards completeness; for instance there is no
section on the automorphisms of (generalized) kummer coverings.

It would be useful to have a charaterization of the Algebra of a
\emph{generalized} Kummer cov.\ (cf.\ with 2.1.8 for Kummer coverings).

\emph{Chap.\ III.} This is a nuisance!

I have tried to write only down those things needed for Chap.\ VIII--X. It was
sufficient to follow the notes which I made for usual schemes and therefore I
have not followed your new set-up. But one way or the other, one gets the
feeling that the notion of tamely ramified covering, although natural, is
very unmanageable if one has to use ``the functor'' instead of the geometric
realization itself.

Lemma 3.2.11 is a technical lemma, used for instance in 3.5.8, 3.6.5 (and
\struck{many} at several places later on).

\emph{Very crucial are} \struck{are} 3.4.1, 3.4.2; also the notations
introduced there are often used in the sequel.

For 3.5.5 I have given only an indication of the proof; in fact many proofs
of III are standard but very unpleasant to spell out in all detail.

Part 3.6 (and especially 3.6.5) are important. Here the geometric
realization determines the tamely ramified covering. Usually only normality,
instead of regularity, seems to be needed. The difficulty is however that the
notion of normality is (perhaps) not stable by étale localization (cf.\ 1.7.3
and 1.7.4).

\emph{Chap.\ VI.} \struck{\ill{}} 6.1.3 is used in 8.1.9, see \add{page} VIII-8
point c. 6.2.1 and 6.2.2 are technical results, used in the proof that 9.3.5
implies 9.3.2 (see \struck{\ill{}} page IX-17).

\page{80}
Comments (2).

\emph{Chap.\ VII.} This is a technical chapter. The results are delicate, but
they are used in an essential way in 9.3 in the reduction step 9.3.3. As
pointed out before: 1.1.13 is used in 7.1.5. In 7.1.5 we have to work with
Kummer coverings instead of generalized Kummer coverings; the reason (or at
least one of the reasons) is that we don't have an analoguous lemma for
2.1.8. As a consequence the assumption $D_i \cap D_j = \emptyset$ (see 9.0.1)
has to be made. Also 7.1.6 is a subtle lemma.

\emph{Chap.\ VIII.} The first page of your notes ``Revêtements formels
étales\ldots''. Remark 8.1.6 seems to me to be correct (!) \add{i.e.\ the
isomorphism is canonical} and essential for 8.1.9.iii.

\emph{Chap.\ IX.} This is a working-out of your notes.

The assumption $D_{i,0} \cap D_{j,0} = \emptyset$ is --- as pointed out above
--- used in 9.3.3, in order to be able to apply 7.1.6 and 7.1.8 (page IX-13,
line 4 top).

\emph{In IX-1} there are several points about which I am not certain (for
instance page IX-4 below). Nevertheless I have used them (!) because it seems
to me that you used them implicitely and without mentioning in your notes and
probably these things are more or less standard results in étale cohomology.

\emph{In IX-2} I have not included (yet) the proof of the key-theorem 9.2.3.
(this proof is explicitly given in the letter of Giraud).

\emph{As to IX-3}: the reduction step 9.3.3. seems to me to be very delicate
and depends upon the results in chap.\ VII. (In your notes you give the
motivation in the margin !). It is --- in particular --- at this point that one
has the feeling that the notion of tamely ramified covering is cumbersome.

I hope that part IX-3 is correctly written \add{and that I don't have
overlooked difficulties}; I had a lot of trouble with it.

One final remark: I needed here the results on tamely ramified covering\add{s}
for \emph{usual schemes}; are these included in your planned final version of
SGA 1 Exp.\ XII (see 9.0.4, 9.1.4, page IX-6 below)?

\emph{Chap.\ X.} The plumbing. X-3 is again very cumbersome. It would be
nicer to work directly with the geometric realization, but there I
encountered difficulties, therefore I followed an easy but long way.

\end{document}
